\documentclass[11pt,a4paper]{article}
\usepackage[margin=23mm]{geometry}
\usepackage{amsmath,amssymb,amsthm}
\setlength{\emergencystretch}{2em}
\newtheorem{theorem}{Theorem}
\newcommand{\gen}[1]{\langle#1\rangle}
\title{The remaining block sizes at a critical head}
\author{Human Compression project: L044 continuation}
\date{}
\begin{document}\maketitle
Retain the complete critical-head hypotheses: $n=2r\ge18$,
$G\le\operatorname{Alt}(\Omega)$ is transitive with connected
string C-marking, and every ambient maximal omission is
intransitive. The acquired head has support $D$ of size six.
Write $E=\Omega\setminus D$, $B=G_{\ge4}$ and
$\sigma=\rho_4$. The group $B$ fixes $D$ pointwise, and
all of its marks other than $\sigma$ fix its anchor $b\in E$.

\begin{theorem}
The group $B$ cannot be transitive on $E$ and preserve a block
system whose block size $a$ and number $z$ are both at least three.
\end{theorem}
\begin{proof}
Every distinguished omission in $B$ is intransitive on $E$:
a transitive omission, together with the head and $\rho_3$
linking it to $b$, would make the corresponding full ambient
omission transitive. Choose a minimal distinguished subset $M$
generating the complete action on the $z$ blocks. Let $c$ count
the actual point orbits of $\gen M$. The acquired actual-orbit
forest and Whiston bounds give
\[
 d-|M|\le c-1,\qquad c\le a,\qquad |M|\le z-1,
 \qquad d=r-4=\frac{az}{2}-1.
\]
Thus $(a-2)(z-2)\le2$. Since $az=|E|\ge12$ is even
and $a,z\ge3$, the only possibilities are $(a,z)=(3,4)$
and $(4,3)$. Both have $|E|=12$, $n=18$, $r=9$ and $d=5$.
In either case every displayed inequality is an equality:
$|M|=z-1$, $c=a$, and the independent block images generate
$S_z$ by Whiston's equality clause. The $a$ actual point
orbits each meet every block once. On each such row the
block action is faithful, and $\gen M$ is the diagonal
natural $S_z$ on the $a$ rows.

For $(a,z)=(3,4)$, an element with odd block image acts
oddly on each of the three rows and hence oddly on $E$.
It fixes $D$, contradicting $\gen M\le G\le A_n$.

For $(a,z)=(4,3)$, the subset $M$ contains $\sigma$:
every other distinguished mark fixes the block containing
$b$, and cannot generate the transitive block action alone.
Its other mark must be $\rho_5$. Indeed the two nonidentity
involutory block images generating $S_3$ are distinct
transpositions and do not commute, whereas $\sigma$ commutes
with every $\rho_j$ for $j\ge6$.
Thus $\gen{\sigma,\rho_5}$ is diagonal natural $S_3$
on four rows of size three.

The nonidentity even mark $\rho_7$ centralizes this entire
group. The centralizer of natural $S_3$ on three points is
trivial, so a permutation centralizing its diagonal action
is exactly a permutation of the four rows, with the unique
equivariant bijection between each pair of rows. As $\rho_7$
fixes $b$, it fixes the row containing $b$ pointwise.
Its remaining action is a nonidentity involution on the
other three rows: the identity row permutation would make
it fix all of $E$ and also $D$, contrary to nonidentity.
It therefore exchanges two rows, making precisely three
point transpositions. This is odd, the final contradiction.
\end{proof}

\paragraph{Attribution and boundary.}
The critical-head data and full actual-orbit forest/Whiston
interfaces are acquired inputs. The new step uses parity in the
forced diagonal actions at equality. Neither a projected IP claim
nor an alternating rank induction, new finite census or pending
packet is used. The cases with block size or number two, primitive
$B$, opposite heads and the unshared endpoint retain their own
proof obligations. No whole-original, formal or question credit
is claimed.
\end{document}
